\thispagestyle{plain}

\section{Introduction}
\label{sec:intro}

Riverine and flash floods recur across northern Bangladesh in most monsoon seasons and cyclone surges strike its south-west coast, yet the susceptibility maps meant to show where flooding would do most harm are rarely tested against floods that actually occurred. Data-driven susceptibility mapping in the country has grown quickly, with a recent review counting 42 distinct conditioning factors in use \citep{islam2025review}, and the studies closest to this one report AUCs between 0.88 and 0.96 \citep{rahman2019, talukdar2020, islamr2024}. Two choices decide what such a figure means, and neither is usually the focus of the paper that reports it. The first is where the flood locations come from, which in most studies is an inventory of a few hundred points drawn from field survey, historical records or a hydrological model. The second is how the test set is separated from the training set, and a random split of spatially clustered points tests a model mostly on locations next to training points, a design known to inflate accuracy \citep{roberts2017, ploton2020}.

Both choices can now be made differently. Sentinel-1 radar sees through monsoon cloud and has been processed into flood maps with overall accuracies above 94\,\% \citep{twele2016, uddin2019}, so every pixel observed to flood can serve as a training label. Whole areas can also be held out when measuring accuracy, so that the test asks how well a model carries to ground it has not seen. A third choice concerns the model itself. Graph neural networks, which pass information between neighbouring locations, have entered hazard mapping mainly for landslides \citep{wang2023, ma2025}, and their appeal for infrastructure is that a road or embankment floods together with its neighbours. Whether that structure adds anything over an ordinary tabular model given the same features has seldom been tested on identical data.

This study addresses three questions, of which the first is the most consequential for how published maps should be read. It asks whether radar-observed flood extents make better training labels than terrain thresholds for ranking flood-prone ground in areas a model has not seen. It then asks whether a graph neural network outperforms ordinary tabular models on the same features, assets and held-out blocks, and finally how far a model trained on earlier floods carries to a later one. The answers are drawn from three flood regions with contrasting regimes, riverine flooding in Rangpur and Rajshahi, flash flooding in the Sylhet haor basin and cyclone surge on the south-west coast, and a landslide model for the Chittagong Hill Tracts completes the hazard coverage. For every cell of a 500\,m grid the flood hazard is combined with exposure, which is what stands on the ground, and vulnerability, which is how hard it would be for the people there to cope, and the results are published through an open dashboard with the data and code behind them.

